\chapter{1999年全国卷（文）}
\section{单选题}

\begin{problem}
如图，\(I\) 是全集，\(M\)，\(P\)，\(S\) 是 \(I\) 的 \(3\text{ 个}\) 子集，则阴影部分所表示的集合是（\quad）

\bitmapfigure[float=inline,width=0.38\linewidth]{img/1999/national_liberal/q1_fig1.png}

\choices
    {\((M\cap P)\cap S\)}
    {\((M\cap P)\cup S\)}
    {\((M\cap P)\cap \overline{S}\)}
    {\((M\cap P)\cup \overline{S}\)}
\end{problem}

\begin{answer}
C
\end{answer}

\begin{solution}
由图可见，阴影部分属于集合 \(M\) 与集合 \(P\) 的公共部分，同时不属于集合 \(S\)，所以阴影部分表示
\(M\cap P\) 与 \(\overline{S}\) 的交集，即 \((M\cap P)\cap\overline{S}\). 因此选 C.
\end{solution}

\begin{problem}
已知映射 \(f: A \to B \)，其中，集合 \(A = \{-3, -2, -1, 1, 2, 3, 4\} \)，集合 \(B\) 中的元素都是 \(A\) 中元素在映射 \(f \) 下的象，且对任意的 \(a \in A \)，在 \(B\) 中和它对应的元素是 \(|a| \)，则集合 \(B\) 中元素的个数是（\quad）

\choices
    {\(4\)}
    {\(5\)}
    {\(6\)}
    {\(7\)}
\end{problem}

\begin{answer}
A
\end{answer}

\begin{solution}
集合 \(B\) 中的元素是集合 \(A\) 中各元素的绝对值. 由
\(\lvert-3\rvert=3\)，\(\lvert-2\rvert=2\)，\(\lvert-1\rvert=1\)，\(\lvert1\rvert=1\)，\(\lvert2\rvert=2\)，\(\lvert3\rvert=3\)，\(\lvert4\rvert=4\)，可知
\(B=\{1,2,3,4\}\). 因而集合 \(B\) 中有 \(4\) 个元素，选 A.
\end{solution}

\begin{problem}
若函数 \(y=f(x)\) 的反函数是 \(y=g(x)\)，\(f(a)=b\)，\(ab\neq0\)，则 \(g(b)\) 等于（\quad）

\choices
    {\(a\)}
    {\(a^{-1}\)}
    {\(b\)}
    {\(b^{-1}\)}
\end{problem}

\begin{answer}
A
\end{answer}

\begin{solution}
由 \(f(a)=b\) 及 \(g=f^{-1}\)，对等式两边施加反函数 \(g\)，得
\(g(f(a))=g(b)\). 因为 \(g(f(a))=a\)，所以 \(g(b)=a\)，选 A.
\end{solution}

\begin{problem}
函数 \(f(x) = M \sin (\omega x + \varphi)\)（\(\omega > 0\)）在区间 \([a, b]\) 上是增函数，且 \(f(a) = -M\)，\(f(b) = M\)，则函数 \(g(x) = M \cos (\omega x + \varphi)\) 在 \([a, b]\) 上（\quad）
\choices
    {是增函数}
    {是减函数}
    {可以取得最大值 \(M\)}
    {可以取得最小值 \(-M\)}
\end{problem}

\begin{answer}
C
\end{answer}

\begin{solution}
记 \(u=\omega x+\varphi\). 由于 \(f\) 在 \([a,b]\) 上递增，且在端点分别取得最小值 \(-M\) 和最大值 \(M\)，相位从正弦函数的谷值沿递增支变化到峰值. 因而可写成
\(u(a)=\frac{3\pi}{2}+2k\pi\)，\(u(b)=\frac{5\pi}{2}+2k\pi\)，其中 \(k\in\Z\)，并且 \(u\) 在该区间内递增.

于是 \(g(x)=M\cos u\) 在此过程中经历从 \(0\) 到 \(M\) 再回到 \(0\) 的变化，故可以取得最大值 \(M\)，选 C.
\end{solution}

\begin{problem}
若 \(f(x) \sin x \) 是周期为 \(\pi \) 的奇函数，则 \(f(x) \) 可以是（\quad）

\choices
    {\(\sin x\)}
    {\(\cos x\)}
    {\(\sin 2x\)}
    {\(\cos 2x\)}
\end{problem}

\begin{answer}
B
\end{answer}

\begin{solution}
逐项检验. 若 \(f(x)=\sin x\)，则 \(f(x)\sin x=\sin^2x\) 是偶函数；若 \(f(x)=\cos x\)，则
\(f(x)\sin x=\sin x\cos x=\frac12\sin2x\)，既是奇函数，周期又为 \(\pi\).

若 \(f(x)=\sin2x\)，则 \(f(x)\sin x=\sin2x\sin x\) 是偶函数；若 \(f(x)=\cos2x\)，则所得函数虽为奇函数，但
\(\sin(x+\pi)\cos2(x+\pi)=-\sin x\cos2x\)，不是周期为 \(\pi\). 因此选 B.
\end{solution}

\begin{problem}
曲线 \(x^{2} + y^{2} + 2\sqrt{2} x - 2\sqrt{2} y = 0\) 关于（\quad）

\choices
    {直线 \(x = \sqrt{2}\) 轴对称}
    {直线 \(y = -x \) 轴对称}
    {点 \((-2, \sqrt{2})\) 中心对称}
    {点 \((- \sqrt{2}, 0)\) 中心对称}
\end{problem}

\begin{answer}
B
\end{answer}

\begin{solution}
配方得
\(x^2+y^2+2\sqrt2x-2\sqrt2y=0\Longleftrightarrow (x+\sqrt2)^2+(y-\sqrt2)^2=4\).
这是圆心为 \((-\sqrt2,\sqrt2)\)、半径为 \(2\) 的圆.

圆关于过圆心的任一直线轴对称. 点 \((-\sqrt2,\sqrt2)\) 在直线 \(y=-x\) 上，因此该圆关于直线 \(y=-x\) 轴对称. 其余选项中的直线不过圆心，或所给点不是圆心，故选 B.
\end{solution}

\begin{problem}
将若干毫升水倒入底面半径为 \(2\mathrm{~cm}\) 的圆柱形器皿中，量得水面的高度为 \(6\mathrm{~cm}\). 若将这些水倒入轴截面是正三角形的倒圆锥形器皿中，则水面的高度是（\quad）

\choices
    {\(6\sqrt{3}\mathrm{~cm}\)}
    {\(6\mathrm{~cm}\)}
    {\(2\sqrt[3]{18}\mathrm{~cm}\)}
    {\(3\sqrt[3]{12}\mathrm{~cm}\)}
\end{problem}

\begin{answer}
B
\end{answer}

\begin{solution}
圆柱内水的体积为
\(V=\pi\cdot2^2\cdot6=24\pi\ \mathrm{cm}^3\).

设倒圆锥中水面的半径为 \(r\)，水面高度为 \(h\). 由于倒圆锥的轴截面为正三角形，圆锥的底面直径与高满足 \(2r=h\sqrt3\)，即 \(r=h/\sqrt3\). 因而水的体积为
\(V=\frac13\pi r^2h=\frac13\pi\cdot\frac{h^2}{3}\cdot h=\frac{\pi h^3}{9}\).

由体积不变，得 \(\frac{\pi h^3}{9}=24\pi\)，所以 \(h^3=216\)，从而 \(h=6\ \mathrm{cm}\). 因此选 B.
\end{solution}

\begin{problem}
若 \((2x + \sqrt{3})^3 = a_0 + a_1x + a_2x^2 + a_3x^3\)，则 \((a_0 + a_2)^2 - (a_1 + a_3)^2\) 的值为（\quad）

\choices
    {\(-1\)}
    {\(1\)}
    {\(0\)}
    {\(2\)}
\end{problem}

\begin{answer}
A
\end{answer}

\begin{solution}
展开左边得
\((2x+\sqrt3)^3=3\sqrt3+18x+12\sqrt3x^2+8x^3\).
所以 \(a_0=3\sqrt3\)，\(a_1=18\)，\(a_2=12\sqrt3\)，\(a_3=8\). 因而
\(a_0+a_2=15\sqrt3\)，\(a_1+a_3=26\)，所求值为
\((15\sqrt3)^2-26^2=675-676=-1\).
因此选 A.
\end{solution}

\begin{problem}
直线 \(\sqrt{3}x+y-2\sqrt{3}=0\) 截圆 \(x^2+y^2=4\) 得的劣弧所对的圆心角为（\quad）

\choices
    {\(\frac{\pi}{6}\)}
    {\(\frac{\pi}{4} \)}
    {\(\frac{\pi}{3} \)}
    {\(\frac{\pi}{2} \)}
\end{problem}

\begin{answer}
C
\end{answer}

\begin{solution}
圆心为原点，半径为 \(2\). 直线到圆心的距离为
\(d=\frac{\lvert-2\sqrt3\rvert}{\sqrt{(\sqrt3)^2+1^2}}=\sqrt3\).

设劣弧所对的圆心角为 \(\theta\)，取弦的一半，则有
\(\cos\frac\theta2=\frac d2=\frac{\sqrt3}{2}\). 因为 \(0<\theta<\pi\)，所以 \(\frac\theta2=\frac\pi6\)，从而 \(\theta=\frac\pi3\). 因此选 C.
\end{solution}

\begin{problem}
如图，在多面体 \(ABCDEF\) 中，已知面 \(ABCD\) 是边长为 \(3\) 的正方形，\(EF\parallel AB\)，\(EF = \frac{3}{2}\)，\(EF\) 与面 \(AC\) 的距离为 \(2\)，则该多面体的体积为（\quad）

\bitmapfigure[float=inline,width=0.38\linewidth]{img/1999/national_liberal/q10_fig1.png}

\choices
    {\(\frac{9}{2}\)}
    {\(5\)}
    {\(6\)}
    {\(\frac{15}{2}\)}
\end{problem}

\begin{answer}
D
\end{answer}

\begin{solution}
由题设距离条件知，上棱 \(EF\) 所在的平面与底面 \(ABCD\) 的距离为 \(2\). 对于底面固定、上棱平行于 \(AB\) 且长度固定的这个多面体，在与底面平行的任一高度处，截面由底面正方形和上棱的相应缩短线段合成；若该高度占总高的比例为 \(t\)，则截面的两边长分别为 \(3(1-t)+\frac32t\) 和 \(3(1-t)\)，所以截面面积与上棱的水平位置无关. 因而可作等积变换，令 \(E\) 在 \(D\) 的正上方，且 \(ED\perp\) 平面 \(ABCD\)，\(ED=2\).

此时该多面体可按公共面 \(FAD\) 分成四棱锥 \(F-ABCD\) 与三棱锥 \(F-ADE\). 四棱锥 \(F-ABCD\) 的高为 \(2\)，其体积为
\(V_{F-ABCD}=\frac13\cdot3^2\cdot2=6\).

又 \(AB\perp AD\)，\(AB\perp ED\)，所以 \(AB\perp\) 平面 \(ADE\)，从而 \(EF\perp\) 平面 \(ADE\). 三角形 \(ADE\) 是直角三角形，且
\(S_{\triangle ADE}=\frac12\cdot AD\cdot ED=\frac12\cdot3\cdot2=3\).
因此三棱锥 \(F-ADE\) 的高为 \(EF=\frac32\)，体积为
\(V_{F-ADE}=\frac13\cdot3\cdot\frac32=\frac32\).

所以多面体的体积为
\(V=6+\frac32=\frac{15}{2}\)，选 D.
\end{solution}

\begin{problem}
若 \(\sin\alpha>\tan\alpha>\cot\alpha\)（\(-\frac{\pi}{2}<\alpha<\frac{\pi}{2}\)），则 \(\alpha\in\)（\quad）

\choices
    {\(\left(-\frac{\pi}{2},-\frac{\pi}{4}\right)\)}
    {\(\left(-\frac{\pi}{4},0\right)\)}
    {\(\left(0,\frac{\pi}{4}\right)\)}
    {\(\left(\frac{\pi}{4},\frac{\pi}{2}\right)\)}
\end{problem}

\begin{answer}
B
\end{answer}

\begin{solution}
在给定区间内 \(\cos\alpha>0\). 由
\(\sin\alpha>\tan\alpha=\frac{\sin\alpha}{\cos\alpha}\)，得
\(\sin\alpha\left(1-\frac1{\cos\alpha}\right)>0\).
由于 \(0<\cos\alpha\leqslant1\)，括号内的因子小于或等于 \(0\)，且不可能取正值，所以必须有 \(\sin\alpha<0\)，即 \(\alpha<0\).

令 \(t=\tan\alpha\). 此时 \(t<0\)，且 \(\tan\alpha>\cot\alpha\) 等价于 \(t>1/t\). 乘以负数 \(t\) 后不等号反向，得 \(t^2<1\)，即 \(-1<t<0\). 因为正切函数在 \((-\frac\pi2,\frac\pi2)\) 上递增，所以
\(-\frac\pi4<\alpha<0\). 因此选 B.
\end{solution}

\begin{problem}
如果圆台的上底面半径为 \(5\)，下底面半径为 \(R\)，中截面把圆台分为上，下两个圆台，它们的侧面积的比为 \(1:2\)，那么 \(R=\)（\quad）

\choices
    {\(10\)}
    {\(15\)}
    {\(20\)}
    {\(25\)}
\end{problem}

\begin{answer}
D
\end{answer}

\begin{solution}
设圆台的母线长为 \(l\). 中截面的半径为 \(\frac{R+5}{2}\)，且上下两个小圆台的母线长均为 \(l/2\). 圆台侧面积与两底面半径之和及母线长的乘积成正比，因此
\[
\frac{S_{\text{上}}}{S_{\text{下}}}
=\frac{5+\frac{R+5}{2}}{\frac{R+5}{2}+R}
=\frac{R+15}{3R+5}=\frac12.
\]
解得 \(2R+30=3R+5\)，所以 \(R=25\). 因此选 D.
\end{solution}

\begin{problem}
给出下列曲线方程：

\begin{circlelist}
\item \(4x+2y-1=0\)；
\item \(x^2+y^2=3\)；
\item \(\frac{x^2}{2}+y^2=1\)；
\item \(\frac{x^2}{2}-y^2=1\).
\end{circlelist}

其中与直线 \(y=-2x-3\) 有交点的所有曲线方程是（\quad）

\choices
    {\(\circled{1}\circled{3}\)}
    {\(\circled{2}\circled{4}\)}
    {\(\circled{1}\circled{2}\circled{3}\)}
    {\(\circled{2}\circled{3}\circled{4}\)}
\end{problem}

\begin{answer}
D
\end{answer}

\begin{solution}
将直线 \(y=-2x-3\) 分别代入四个曲线方程检验交点.

对于 \(\circled{1}\)，得 \(4x+2(-2x-3)-1=-7\neq0\)，所以没有交点.

对于 \(\circled{2}\)，代入得
\(x^2+(-2x-3)^2=3\)，即 \(5x^2+12x+6=0\). 其判别式为 \(12^2-4\cdot5\cdot6=24>0\)，所以有交点.

对于 \(\circled{3}\)，代入得
\(\frac{x^2}{2}+(-2x-3)^2=1\)，即 \(9x^2+24x+16=0\)，也就是 \((3x+4)^2=0\)，所以有交点.

对于 \(\circled{4}\)，代入得
\(\frac{x^2}{2}-(-2x-3)^2=1\)，即 \(7x^2+24x+20=0\). 其判别式为 \(24^2-4\cdot7\cdot20=16>0\)，所以有交点.

因此有交点的曲线是\(\circled{2}\)、\(\circled{3}\)、\(\circled{4}\)，选 D.
\end{solution}

\begin{problem}
某电脑用户计划使用不超过 \(500\text{ 元}\) 的资金购买单价分别为 \(60\text{ 元}\)，\(70\text{ 元}\) 的单片软件和盒装磁盘，根据需要，软件至少买 \(3\text{ 片}\)，磁盘至少买 \(2\text{ 盒}\)，则不同的选购方式共有（\quad）

\choices
    {\(5\text{ 种}\)}
    {\(6\text{ 种}\)}
    {\(7\text{ 种}\)}
    {\(8\text{ 种}\)}
\end{problem}

\begin{answer}
C
\end{answer}

\begin{solution}
设购买软件 \(x\) 片、磁盘 \(y\) 盒，则 \(x,y\in\Z\)，且
\(x\geqslant3\)，\(y\geqslant2\)，\(60x+70y\leqslant500\).

固定 \(x\) 计数. 当 \(x=3\) 时，\(140\leqslant70y\leqslant320\)，所以 \(y=2,3,4\)，有 \(3\) 种；当 \(x=4\) 时，\(140\leqslant70y\leqslant260\)，所以 \(y=2,3\)，有 \(2\) 种；当 \(x=5\) 时，\(y=2\)，有 \(1\) 种；当 \(x=6\) 时，\(y=2\)，有 \(1\) 种；当 \(x\geqslant7\) 时，即使 \(y=2\) 也超过 \(500\) 元.

所以不同选购方式共有 \(3+2+1+1=7\) 种，选 C.
\end{solution}

\section{填空题}

\begin{problem}
设椭圆 \(\frac{x^2}{a^2}+\frac{y^2}{b^2}=1\)（\(a>b>0\)）的右焦点为 \(F_1\)，右准线为 \(l_1\)，若过 \(F_1\) 且垂直于 \(x\) 轴的弦长等于点 \(F_1\) 到 \(l_1\) 的距离，则椭圆的离心率是\fillinblank{}.
\end{problem}

\begin{answer}
\(\frac{1}{2}\)
\end{answer}

\begin{solution}
设椭圆半焦距为 \(c\)，则 \(c^2=a^2-b^2\)，离心率为 \(e=c/a\). 过右焦点的垂直于 \(x\) 轴的弦所在直线为 \(x=c\). 代入椭圆方程，弦长为
\(2b\sqrt{1-\frac{c^2}{a^2}}=\frac{2b^2}{a}\).

右准线方程为 \(x=\frac{a^2}{c}\)，所以点 \(F_1\) 到右准线的距离为
\(\frac{a^2}{c}-c=\frac{a^2-c^2}{c}=\frac{b^2}{c}\).

由题意 \(\frac{2b^2}{a}=\frac{b^2}{c}\). 因为 \(b>0\)，可约去 \(b^2\)，得 \(2c=a\)，所以 \(e=\frac ca=\frac12\).
\end{solution}

\begin{problem}
在一块并排 \(10\text{ 垄}\) 的田地中，选择 \(2\text{ 垄}\) 分别种植 \(A\)，\(B\) 两种作物，每种作物种植一垄，为有利于作物生长，要求 \(A\)，\(B\) 两种作物的间隔不小于 \(6\text{ 垄}\)，则不同的选垄方法共有\fillinblank{}\(\text{种}\).
\end{problem}

\begin{answer}
12
\end{answer}

\begin{solution}
设种植 \(A\) 和 \(B\) 的垄号分别为 \(i\) 和 \(j\)，不妨设 \(i<j\). “间隔不少于 \(6\) 垄”表示两垄之间至少有 \(6\) 垄，因此 \(j-i\geqslant7\). 可行的无序垄号对为
\((1,8),(1,9),(1,10),(2,9),(2,10),(3,10)\)，共 \(6\) 对.

每一对垄号中，\(A\)，\(B\) 的种植顺序有 \(2\) 种，所以不同的选垄方法共有 \(6\times2=12\) 种.
\end{solution}

\begin{problem}
若正数 \(a\)，\(b\) 满足 \(ab=a+b+3\)，则 \(ab\) 的取值范围是\fillinblank{}.
\end{problem}

\begin{answer}
\([9,+\infty)\)
\end{answer}

\begin{solution}
令 \(p=ab\). 由题设 \(a+b=p-3\). 因为 \(a,b\) 是方程
\(u^2-(p-3)u+p=0\) 的两个正根，所以其判别式必须满足
\((p-3)^2-4p\geqslant0\)，即 \((p-1)(p-9)\geqslant0\).

又因 \(a,b>0\)，有 \(p=ab=a+b+3>3\)，所以不可能有 \(p\leqslant1\)，从而 \(p\geqslant9\). 反之，对任意 \(p\geqslant9\)，上述方程的两根和 \(p-3>0\)、积 \(p>0\)，且判别式非负，故两根均为正数，能够满足题设. 因此 \(ab\in[9,+\infty)\).
\end{solution}

\begin{problem}
\(\alpha\)，\(\beta\) 是两个不同的平面，\(m\)，\(n\) 是平面 \(\alpha\) 及 \(\beta\) 之外的两条不同直线. 给出四个论断：

\begin{circlelist}
\item \(m \perp n\)；
\item \(\alpha \perp \beta\)；
\item \(n \perp \beta\)；
\item \(m \perp \alpha\).
\end{circlelist}

以其中三个论断作为条件，余下一个论断作为结论，写出你认为正确的一个命题：\fillinblank{}.
\end{problem}

\begin{answer}
\(m\perp\alpha\)，\(n\perp\beta\)，\(\alpha\perp\beta\) \(\Longrightarrow\) \(m\perp n\)
\end{answer}

\begin{solution}
由 \(\alpha\perp\beta\) 且 \(m\perp\alpha\)，可知直线 \(m\) 的方向垂直于平面 \(\alpha\)，因而平行于平面 \(\beta\). 又因 \(n\perp\beta\)，直线 \(n\) 垂直于平面 \(\beta\) 内任一直线及其平行线，所以 \(n\perp m\). 因此所写命题正确.
\end{solution}

\section{解答题}

\begin{problem}
解方程：\(\sqrt{3\lg x-2}-3\lg x+4=0\).
\end{problem}

\begin{solution}
由根式有意义，得 \(x>0\) 且 \(3\lg x-2\geqslant0\). 令 \(t=\sqrt{3\lg x-2}\)，则 \(t\geqslant0\)，且 \(3\lg x-2=t^2\). 原方程化为
\[
t-t^2+2=0\quad\Longleftrightarrow\quad (t-2)(t+1)=0.
\]
由于 \(t\geqslant0\)，所以 \(t=2\)，从而 \(3\lg x-2=4\)，即 \(\lg x=2\)，得 \(x=100\). 经检验，\(x=100\) 满足原方程，因此原方程的解为 \(x=100\).
\end{solution}

\begin{problem}
数列 \(\{a_{n}\}\) 的前 \(n\) 项和记为 \(S_{n}\). 已知 \(a_{n} = 5S_{n} - 3\)（\(n \in \N\)），求 \(\displaystyle\lim_{n \to \infty}(a_{1} + a_{3} + \cdots + a_{2n-1})\) 的值.
\end{problem}

\begin{solution}
当 \(n\geqslant2\) 时，由 \(a_n=S_n-S_{n-1}\) 及已知条件，得
\[
a_n-a_{n-1}=5(S_n-S_{n-1})=5a_n,
\]
所以 \(a_n=-\frac14a_{n-1}\). 当 \(n=1\) 时，\(S_1=a_1\)，于是
\[
a_1=5S_1-3=5a_1-3,
\]
解得 \(a_1=\frac34\). 因此数列 \(\{a_n\}\) 是首项为 \(\frac34\)、公比为 \(-\frac14\) 的等比数列.

于是奇数项组成首项为 \(a_1=\frac34\)、公比为 \(\left(-\frac14\right)^2=\frac1{16}\) 的等比数列，故
\[
a_1+a_3+\cdots+a_{2n-1}=\frac34\cdot\frac{1-\left(\frac1{16}\right)^n}{1-\frac1{16}}.
\]
所以
\[
\lim_{n\to\infty}(a_1+a_3+\cdots+a_{2n-1})
=\frac34\cdot\frac1{1-\frac1{16}}=\frac45.
\]
\end{solution}

\begin{problem}
设复数 \(z=3\cos\theta+\i\cdot 2\sin\theta\)，求函数 \(y=\theta-\arg z\)（\(0<\theta<\frac{\pi}{2}\)）的最大值以及对应的 \(\theta\) 值.
\end{problem}

\begin{solution}
由 \(0<\theta<\frac\pi2\) 知 \(z\) 在第一象限，设 \(\alpha=\arg z\)，则 \(0<\alpha<\frac\pi2\)，且
\[
\tan\alpha=\frac{2\sin\theta}{3\cos\theta}=\frac23\tan\theta.
\]
因此 \(0<y=\theta-\alpha<\frac\pi2\). 令 \(t=\tan\theta>0\)，利用正切差角公式，得
\[
\tan y=\frac{t-\frac23t}{1+\frac23t^2}=\frac{t}{3+2t^2}.
\]
又因为
\[
\frac{3+2t^2}{t}=\frac3t+2t\geqslant2\sqrt6,
\]
所以
\[
0<\tan y=\frac{t}{3+2t^2}\leqslant\frac1{2\sqrt6}=\frac{\sqrt6}{12}.
\]
等号成立当且仅当 \(\frac3t=2t\)，即 \(t=\sqrt{\frac32}=\frac{\sqrt6}{2}\). 由于正切函数在 \(\left(0,\frac\pi2\right)\) 上单调递增，故函数 \(y\) 的最大值为
\[
\arctan\frac{\sqrt6}{12},
\]
对应的 \(\theta=\arctan\frac{\sqrt6}{2}\).
\end{solution}

\begin{problem}
如图，已知正四棱柱 \(ABCD-A_1B_1C_1D_1\)，点 \(E\) 在棱 \(D_1D\) 上，截面 \(EAC\parallel D_1B\)，且面 \(EAC\) 与底面 \(ABCD\) 所成的角为 \(45^\circ\)，\(AB=a\).

\begin{enumerate}
\item 求截面 \(EAC\) 的面积；
\item 求异面直线 \(A_1B_1\) 与 \(AC\) 之间的距离；
\item 求三棱锥 \(B_1-EAC\) 的体积.
\end{enumerate}

\bitmapfigure[float=inline,width=0.38\linewidth]{img/1999/national_liberal/q22_fig1.png}
\end{problem}

\begin{solution}
\begin{enumerate}
\item 设 \(O=AC\cap BD\)，连接 \(EO\). 因为 \(ABCD\) 是正方形，所以 \(DO\perp AC\). 又 \(ED\perp\) 底面 \(ABCD\)，所以 \(ED\perp AC\). 直线 \(ED\) 与 \(DO\) 是平面 \(EDO\) 内相交的两条直线，且都垂直 \(AC\)，故 \(AC\perp\) 平面 \(EDO\)，从而 \(EO\perp AC\). 因此 \(\angle EOD\) 是平面 \(EAC\) 与底面 \(ABCD\) 所成二面角的平面角，故 \(\angle EOD=45^\circ\). 在直角三角形 \(EOD\) 中，\(DO=\frac{\sqrt2}{2}a\)，所以
\[
EO=DO\sec45^\circ=a.
\]
又 \(AC=\sqrt2a\)，故
\[
S_{\triangle EAC}=\frac12AC\cdot EO=\frac{\sqrt2}{2}a^2.
\]

\item 由于 \(A_1A\perp\) 底面 \(ABCD\)，所以 \(A_1A\perp AC\)，又 \(A_1A\perp A_1B_1\)，故 \(A_1A\) 是异面直线 \(A_1B_1\) 与 \(AC\) 的公垂线.

在平面 \(D_1DB\) 内，平面 \(EAC\) 与它的交线为 \(EO\). 由 \(D_1B\parallel\) 平面 \(EAC\)，得 \(D_1B\parallel EO\). 在三角形 \(D_1DB\) 中，\(O\) 是 \(DB\) 的中点，且 \(EO\parallel D_1B\)，所以 \(E\) 是 \(D_1D\) 的中点. 由中位线定理，\(D_1B=2EO=2a\). 又 \(DB=\sqrt2a\)，且 \(D_1D\perp DB\)，于是
\[
D_1D=\sqrt{D_1B^2-DB^2}=\sqrt{4a^2-2a^2}=\sqrt2a.
\]
正四棱柱的高 \(A_1A=D_1D\)，故异面直线 \(A_1B_1\) 与 \(AC\) 之间的距离为 \(\sqrt2a\).

\item 取 \(A\) 为原点，以 \(AB\)，\(AD\)，\(AA_1\) 所在直线为坐标轴. 由上问知 \(AA_1=\sqrt2a\)，且 \(E\) 是 \(D_1D\) 的中点，所以
\(A(0,0,0)\)，\(C(a,a,0)\)，\(E\left(0,a,\frac{a}{\sqrt2}\right)\)，\(B_1(a,0,\sqrt2a)\).
由三点 \(A\)，\(C\)，\(E\) 得平面 \(EAC\) 的方程为
\[
x-y+\sqrt2z=0.
\]
点 \(B_1\) 到该平面的距离为
\[
h=\frac{\left|a-0+\sqrt2\cdot\sqrt2a\right|}{\sqrt{1+1+2}}=\frac{3a}{2}.
\]
因此
\[
V_{B_1-EAC}=\frac13S_{\triangle EAC}\cdot h
=\frac13\cdot\frac{\sqrt2}{2}a^2\cdot\frac{3a}{2}=\frac{\sqrt2}{4}a^3.
\]
\end{enumerate}
\end{solution}

\begin{problem}
如图为一台冷轧机的示意图. 冷轧机由若干对轧辊组成，带钢从一端输入，经过各对轧辊逐步减薄后输出.

\begin{enumerate}
\item 输入带钢的厚度为 \(\alpha\)，输出带钢的厚度为 \(\beta\)，若每对轧辊的减薄率不超过 \(r_0\). 问冷轧机至少需要安装多少对轧辊？
\item 已知一台冷轧机共有 \(4\) 对减薄率为 \(20\%\) 的轧辊，所有轧辊周长均为 \(1600\mathrm{~mm}\). 若第 \(k\) 对轧辊有缺陷，每滚动一周在带钢上压出一个疵点，在冷轧机输出的带钢上，疵点的间距为 \(L_k\). 为了便于检修，请计算 \(L_1\)，\(L_2\)，\(L_3\) 并填入下表（轧钢过程中，带钢宽度不变，且不考虑损耗）.
\end{enumerate}

（一对轧辊减薄率 \(=\frac{\text{输入该对的带钢厚度}-\text{从该对输出的带钢厚度}}{\text{输入该对的带钢厚度}}\)）

\begin{center}
\begin{tabular}{|c|c|c|c|c|}
\hline
轧辊序号 & \(1\) & \(2\) & \(3\) & \(4\) \\
\hline
疵点间距 \(L_k/\mathrm{mm}\) &  &  &  & \(1600\) \\
\hline
\end{tabular}
\end{center}

\FigureLayoutDeclare{img/1999/national_liberal/q23_fig1.png}{7.5cm}{38bp 44bp 27bp 34bp}
\bitmapfigure[float=inline,width=0.95\linewidth]{img/1999/national_liberal/q23_fig1.png}
\end{problem}

\begin{solution}
\begin{enumerate}
\item 设安装 \(n\) 对轧辊后带钢的厚度为 \(t_n\). 每对轧辊的减薄率不超过 \(r_0\)，所以每经过一对轧辊，厚度至少保留原来的 \(1-r_0\)，从而
\[
t_n\geqslant\alpha(1-r_0)^n.
\]
要使带钢厚度能够减薄到不超过 \(\beta\)，必须有
\[
\alpha(1-r_0)^n\leqslant\beta.
\]
由于 \(0<1-r_0<1\)，两边取对数并注意 \(\lg(1-r_0)<0\)，得
\[
n\geqslant\frac{\lg\beta-\lg\alpha}{\lg(1-r_0)}.
\]
因此至少需要安装
\[
n_{\min}=\left\lceil\frac{\lg\beta-\lg\alpha}{\lg(1-r_0)}\right\rceil
\]
对轧辊.

\item 设带钢宽度为 \(w\)，第 \(k\) 对轧辊出口处的厚度为 \(t_k=\alpha(0.8)^k\). 第 \(k\) 对轧辊压出的两个相邻疵点之间的带钢长度为 \(1600\mathrm{~mm}\)，其体积为 \(1600t_kw\). 经过第 \(4\) 对轧辊后，这段带钢的长度为 \(L_k\)，厚度为 \(t_4=\alpha(0.8)^4\)，体积为 \(L_kt_4w\). 因为轧钢过程中不考虑损耗，且带钢宽度不变，所以
\[
1600t_kw=L_kt_4w,
\]
即
\[
L_k=1600\frac{t_k}{t_4}=1600\cdot(0.8)^{k-4}.
\]
因此 \(L_1=3125\mathrm{~mm}\)，\(L_2=2500\mathrm{~mm}\)，\(L_3=2000\mathrm{~mm}\).
\end{enumerate}
\end{solution}

\begin{problem}
如图，给出定点 \(A(a,0)\)（\(a>0\)）和直线 \(l\colon x=-1\). \(B\) 是直线 \(l\) 上的动点，\(\angle BOA\) 的角平分线交 \(AB\) 于点 \(C\). 求点 \(C\) 的轨迹方程，并讨论方程表示的曲线类型与 \(a\) 值的关系.

\bitmapfigure[float=inline,width=0.33\linewidth]{img/1999/national_liberal/q24_fig1.png}
\end{problem}

\begin{solution}
先设 \(B=(-1,t)\)，其中 \(t\in\R\)，并记 \(OB=r=\sqrt{1+t^2}\). 当 \(t\neq0\) 时，在三角形 \(AOB\) 中，由角平分线定理
\[
\frac{AC}{CB}=\frac{AO}{OB}=\frac ar.
\]
因此点 \(C\) 的坐标为内分点坐标
\[
C\left(\frac{r\cdot a+a\cdot(-1)}{a+r},\frac{a\cdot t}{a+r}\right)
=\left(\frac{a(r-1)}{a+r},\frac{at}{a+r}\right).
\]
记 \(C(x,y)\)，则
\(x=\frac{a(r-1)}{a+r}\)，\(y=\frac{at}{a+r}\).
由于 \(r>1\)，可知 \(0<x<a\). 又 \(t^2=r^2-1\)，代入上式可得
\[
\begin{aligned}
(a+1)y^2
&=\frac{a^2(a+1)(r^2-1)}{(a+r)^2}\\
&=\frac{a^2\bigl((a-1)(r-1)^2+2(r-1)(a+r)\bigr)}{(a+r)^2}\\
&=(a-1)x^2+2ax.
\end{aligned}
\]
所以点 \(C\) 的轨迹方程为
\[
(1-a)x^2-2ax+(1+a)y^2=0.
\]
当 \(t=0\) 时，\(B=(-1,0)\)，此时 \(AB\) 与 \(OA\) 共线，角平分线与 \(AB\) 交于 \(O\)，仍得到 \(C=(0,0)\). 因此轨迹方程应结合条件 \(0\leqslant x<a\) 取值.

反过来，对满足该方程及 \(0<x<a\) 的点，令 \(r=\dfrac{a(1+x)}{a-x}\)，再由 \(t=\dfrac{y(a+r)}a\) 定义 \(B=(-1,t)\)，可验证 \(r^2=1+t^2\)，故这些点均可由相应的动点 \(B\) 得到；当 \(x=0\) 时，方程给出 \(y=0\)，正对应于上面的 \(B=(-1,0)\). 所以轨迹为上述方程在 \(0\leqslant x<a\) 内的部分.

当 \(0<a<1\) 时，方程可化为
\[
(1-a)\left(x-\frac{a}{1-a}\right)^2+(1+a)y^2=\frac{a^2}{1-a},
\]
表示椭圆的一段.

当 \(a=1\) 时，方程化为 \(y^2=x\)，表示抛物线的一段.

当 \(a>1\) 时，方程可化为
\[
(a-1)\left(x+\frac{a}{a-1}\right)^2-(a+1)y^2=\frac{a^2}{a-1},
\]
表示双曲线的一段.
\end{solution}
